\documentclass[12pt]{article}
\usepackage[utf8]{inputenc}
\usepackage[a4paper, margin=2cm]{geometry}
\usepackage{amsthm}
\usepackage{amsmath, amssymb}
\usepackage{circuitikz}
\usepackage{graphicx}
\usepackage{karnaugh-map}
\usepackage{newtxtext, newtxmath}
\usepackage{tcolorbox}
\usepackage{tikz}
\usepackage{titlesec}
\usepackage{xcolor}
\usepackage{xeCJK}
\usepackage{listings}
\definecolor{lightblue}{RGB}{135, 206, 250}
\tcbuselibrary{skins,theorems}

\titleformat{\section}[block]{\Large\bfseries\centering}{}{0pt}{}
\titleformat{\subsection}[block]{\bfseries}{\thesubsection}{2pt}{}

\newtcbtheorem[no counter]{proposition}{\large\hspace{-8pt}Proposition}{
  enhanced,
  colback=white,
  colframe=lightblue,
  left=12pt,
  right=12pt,
  top=24pt,
  bottom=24pt,
  fonttitle=\bfseries
}{prop}
\newtcbtheorem[]{question}{\large\hspace{-8pt}Question :}{
  enhanced,
  colback=white,
  colframe=black,
  left=12pt,
  right=12pt,
  top=24pt,
  bottom=24pt,
  fonttitle=\bfseries
}{qst}
\newenvironment{Answer}{
  \vspace{12pt}
  {\large\textbf{\textit{\hspace{-12pt}Answer:}}}\par
  \vspace{8pt}
} {\vspace{24pt}}
\newenvironment{Proof}{
  \vspace{12pt}
  {\large\textbf{\textit{\hspace{-12pt}Proof:}}}\par
  \vspace{8pt}
} {\vspace{24pt}}

\lstset{
    language=Python,
    basicstyle=\ttfamily\footnotesize,
    keywordstyle=\color{blue}\bfseries,
    commentstyle=\color{gray},
    stringstyle=\color{red},
    showstringspaces=false,
    numberstyle=\tiny\color{gray},
    stepnumber=1,
    numbersep=10pt,
    frame=single,
    breaklines=true,
    captionpos=b
}

\title{\textbf{Algorithm Design Homework 3}}
\date{June 18, 2024}

\begin{document}
\maketitle

\begin{question}{}{}
找零钱问题：假设一个硬币系统有三种硬币(1分，5分，10分)，现在要找给顾客一角五分的零钱。每次拿出不超过零钱数量的最大面值硬币 (即贪心选择策略)，将使所拿出的硬币个数最少。
\begin{itemize}
  \item (1)证明:硬币系统 (1分，5分，10分)通过贪心选择策略能够得到任意零钱n的最优解。
  \item (2)分析:为什么硬币系统 (1分，5分，6分) 通过贪心选择策略不一定能够得到任意零钱n的最优解。
  \item (3)设计: 硬币系统 (1分，5分，6分)要想得到任意零钱n的最优解，如何设计算法?
\end{itemize}
\end{question}
\begin{Answer}
\subsection*{(1)贪心策略的证明}
\noindent

\textbf{贪心选择策略}指的是每次尽可能选择面值最大的硬币，直到达到所需的总金额。

对于硬币系统 (1分，5分，10分)：

\begin{itemize}
  \item 假设我们需要找 n 分的零钱。
  \item 贪心策略的步骤是：每次尽量使用 10 分的硬币，然后是 5 分的硬币，最后是 1 分的硬币。
\end{itemize}

我们来证明这种策略总是能得到最优解。

证明：
假设我们要找 n 分的零钱，用贪心策略得到的硬币数是最优的。
如果我们找到的硬币组合不是最优的，那应该存在另一种组合，硬币数量更少。

\begin{itemize}
  \item 对于任意 n，当 n ≥ 10 时，使用一个 10 分硬币后，剩余的钱仍然可以用相同的策略处理。
  \item 如果 10 ≤ n < 15，可以用一个 10 分硬币加上剩余的 1-4 分，这样总是最少的硬币数。
  \item 如果 5 ≤ n < 10，可以用一个 5 分硬币加上剩余的 1-4 分，这样也是最少的硬币数。
  \item 如果 1 ≤ n < 5，可以直接用 1 分硬币，这也是最少的硬币数。
\end{itemize}

因此，贪心策略在每一步都是最优的，最终得到的解也是最优解。

\subsection*{(2)贪心策略不可行的分析}
\noindent
对于硬币系统 (1分，5分，6分)：

\begin{itemize}
  \item 假设我们需要找 n 分的零钱。
  \item 贪心策略的步骤是：每次尽量使用 6 分的硬币，然后是 5 分的硬币，最后是 1 分的硬币。
\end{itemize}

我们来看一个反例：

假设需要找 10 分的零钱：

\begin{itemize}
  \item 贪心策略会首先选择一个 6 分的硬币，然后还需要找 4 分的零钱。
    这时只能用 1 分的硬币找到 4 分，总共需要 5 个硬币（1个6分 + 4个1分）。
  \item 但是，最优解应该是使用两个 5 分硬币，总共需要 2 个硬币。
\end{itemize}

这个例子说明了贪心策略不能保证总是找到最优解，因为选择最大面值硬币并不一定总是能得到硬币数量最少的组合。

\subsection*{(3)算法设计}

为了解决这个问题，我们可以使用动态规划（Dynamic Programming）算法来找到最优解。

动态规划算法设计：
我们将使用一个二维数组 dp，其中 dp[i][j] 表示找到 i 分零钱时，用硬币 j 的最少数量。
我们还将使用一个数组 coin\_count 来记录每个金额所需的硬币种类的数量。

\begin{itemize}
  \item 初始化一个数组 dp，大小为 (amount + 1) x len(coins)，初始值为 0，表示每种硬币的数量。
  \item 初始化一个数组 min\_coins，大小为 amount + 1，初始值为无穷大，表示找到每个金额所需的最少硬币数。min\_coins[0] = 0 表示 0 分零钱需要 0 个硬币。
  \item 对于每个金额 i，从 1 到 amount，遍历每种硬币面值 coin，如果当前面值小于等于 i，更新 min\_coins[i] 和 dp[i] 为最优值。
\end{itemize}

下面是详细的算法实现：

\begin{lstlisting}
def min_coins(coins, amount):
    # Using amount + 1 represents the infinite value.
    dp = [[0] * len(coins) for _ in range(amount + 1)]
    min_coins = [amount + 1] * (amount + 1)
    min_coins[0] = 0

    for i in range(1, amount + 1):
        for j, coin in enumerate(coins):
            if i - coin >= 0:
                if min_coins[i - coin] + 1 < min_coins[i]:
                    min_coins[i] = min_coins[i - coin] + 1
                    dp[i] = dp[i - coin].copy()
                    dp[i][j] += 1

    if min_coins[amount] == amount + 1:
        return -1, []
    else:
        return min_coins[amount], dp[amount]
\end{lstlisting}

解释：

\begin{itemize}
  \item dp 数组用于记录每种硬币的数量。
  \item min\_coins 数组用于记录每个金额所需的最少硬币数。
  \item 在更新 min\_coins[i] 时，如果发现使用当前硬币 coin 能减少硬币数量，就更新 dp[i] 数组中对应的硬币数量。
\end{itemize}

输出：

\begin{itemize}
  \item 最少硬币数。
  \item 每种硬币的数量。
\end{itemize}

这样我们不仅得到了最少硬币数，同时也得到了每种硬币的具体数量。

这个算法的时间复杂度是$O(n×k)$，其中 n 是金额，k 是硬币种类数。
通过动态规划，我们可以确保找到每个金额的最优解。

\end{Answer}

\begin{flushright}\begin{tabular}{rl}
  姓名: & 李俊洁\\
  班级: & 计算机2204 \\
  学号: & 20225443 \\
  课程: & 算法设计与分析 \\
  日期: & June 18, 2024 \\
  制作: & \LaTeX \\
  & 感谢您的批阅！
\end{tabular}\end{flushright}
\end{document}
